\begin{abstract}
Electric easy-bikes carry much of Bangladesh's short-distance transport, yet most run on lead-acid batteries that last about \lifeLA{} months and feed an informal recycling chain that is a leading source of lead exposure. Whether lithium iron phosphate (LFP) packs should replace them depends on field cost and life, not on laboratory cycle life alone. We combine a field survey of \nSurvey{} easy-bike drivers (October 2026) with NASA POWER temperature for Khulna, Chuadanga and Rajshahi (2024--2025) in a MATLAB/Simulink model of capacity, resistance, depth of discharge, Arrhenius ageing, replacement-aware cost, carbon and lead flow. Survey prices (median BDT \costMedLA{} lead-acid, \costMedLFP{} LFP), replacement intervals (\replMedLA{} and \replMedLFP{} months), resale values and daily distance (\kmAll{} km) replace literature assumptions, and the model is checked against the survey and published energy per charge. Literature cycle lives would make LFP BDT \litGap{} cheaper over ten years, but survey-anchored lives make it BDT \gapLFP{} (\gapPct\%) dearer for the operator. LFP breaks even at a set price of about BDT \beP{} or a life of about \beL{} months. It still cuts grid electricity by \gridratio\%, saves about \saveBike{} kg CO$_2$e per bike-year and removes \leadLow--\leadHigh{} kt of lead per year from recycling at 1.4 million bikes, at an implied BDT \abateCarbon{} per tonne CO$_2$e. Mean battery temperatures at the three sites differ by only \Tspread\,$^\circ$C, so site does not change the ranking. The survey is a small convenience sample and the LFP results rest on five drivers.
\end{abstract}
\noindent\textbf{Keywords:} easy-bike; lead-acid battery; lithium iron phosphate; field survey; NASA POWER; battery degradation; total cost of ownership; carbon emissions; lead pollution; Bangladesh
\bigskip
